\section{Off-Belief Learning}
\label{sec:method}
One of the big challenges in ZSC under partially observable settings is to determine how to \emph{interpret} the actions of other agents and how to select actions that will be interpretable to other agents. 
In the following we introduce OBL, a method that addresses this problem by learning a hierarchy of policies, with an optimal \emph{grounded} policy at the lowest level, which does not interpret other agents' actions at all.
We will discuss OBL both as a formal method with corresponding proofs and as a scalable algorithm based on a \emph{fictitious transition} mechanism applicable to the deep RL setting.

\subsection{The OBL Operator}
We first introduce the OBL operator that computes $\pi_1$ given any $\pi_0$. 
If a common knowledge policy $\pi_0$ is played by all agents up to $\tau^i$, then agent $i$ can compute a belief distribution $\mathcal{B}_{\pi_0}(\tau|\tau^i)=P(\tau|\tau^i,\pi_0)$ 
conditional on its AOH. This belief distribution fully characterizes the effect of the history on the current state. Notice that the return for (all players) playing a `counterfactual' policy $\pi_0$ 
to $\tau^i$ and $\pi_1$ thereafter, which we denote $V^{\pi_0 \to \pi_1}(\tau^i)$, is precisely the expected return for all players playing according to $\pi_1$, starting from a trajectory $\tau \sim \mathcal{B}_{\pi_0}(\tau^i)$. Therefore we define a counterfactual value function as follows:
\begin{equation}
    V^{\pi_0 \to \pi_1}(\tau^i) = \mathbb{E}_{\tau \sim \mathcal{B}_{\pi_0}(\tau^i)}\left[ V^{\pi_1}(\tau) \right].
    \label{eq:V_counterfactual}
\end{equation}
We can similarly define counterfactual $Q$ values as
\begin{multline}
    Q^{\pi_0 \to \pi_1}(a|\tau_t^i) = \sum_{\tau_t,\tau_{t+1}}\mathcal{B}_{\pi_0}(\tau_t|\tau_t^i) \big[ R(s_t, a) \\
    \hspace{2cm} + \mathcal{T}(\tau_{t+1}|\tau_t) V^{\pi_1}(\tau_{t+1})\big].
    \label{eq:Q_counterfactual}
\end{multline}

Then \textit{OBL} operator is defined to be the operator that maps an initial policy $\pi_0$ to a new policy $\pi_1$ as follows:

\begin{equation}
    \label{eq:pi1}
    \pi_1(a|\tau^i)=\frac{\exp (Q^{\pi_0\to \pi_1}(a|\tau^i)/T) }{\sum_{a'} \exp (Q^{\pi_0\to \pi_1}(a'|\tau^i)/T)}
\end{equation}
for some temperature hyperparameter $T$.

\subsection{Properties of Off-Belief Learning}

\begin{restatable}{theorem}{thmunique}
\label{thm:unique}
For any $T>0$ and starting policy $\pi_0$, OBL computes a unique policy $\pi_1$.
\end{restatable}
Proof in the Appendix~\ref{app:proofs}. ``Unique'' means that that\sam{,} once $T$ is fixed, $\pi_1$ is a well-defined deterministic function of $\pi_0$. This means that, unlike most MARL methods, OBL always converges to the same answer in principle, regardless of random initialization or other hyperparameters.
\vspace{0.15cm}
\begin{corollary}
For any $T>0$ and starting policy $\pi_0$, $N$ agents independently computing OBL joint policy $\pi_1 = \{\pi_1^i\}$ and playing their parts of it ($\pi_1^i$) will achieve the same return as if they had computed a centralized $\pi_1$ policy (zero-shot coordination).
\label{cor:pi1}
\end{corollary}
\begin{proof}
Since all $\pi_1$ are identical, this follows trivially.
\end{proof}

Theorem $\ref{thm:unique}$ is trivial in a single-agent context, but illustrates a substantial departure from traditional multi-agent learning rules, under which independently computed policies for each agent are typically \textit{not} unique or compatible due to the presence of multiple equilibria, a particularly severe example of which is the formation of `arbitrary' conventions for communication in a game like Hanabi.

\begin{restatable}{theorem}{thmpolicyimprovement}
For every policy $\pi_1$ generated by OBL from $\pi_0$, $J(\pi_1) \geq J(\pi_0) - t_{max}T/e$, i.e. OBL is a policy improvement operator except for a term that vanishes as $T \to 0$.
\label{thm:policy_improvement}
\end{restatable}

Proof in the Appendix~\ref{app:proofs}. Theorem~\ref{thm:policy_improvement} implies that self-play performance will gradually increase if we apply OBL iteratively. Together with Corollary~\ref{cor:pi1}, it also guarantees that zero-shot coordination performance, i.e. cross-play between independent training runs, will improve at the same pace.

Notably, unlike standard MARL, the fixed points of the OBL learning rule \textit{are not} guaranteed to be equilibria of the game.
However, we have the theorem:
\begin{restatable}{theorem}{thmequilibrium}
    If repeated application of the OBL policy improvement operator converges to a fixed point policy $\pi$, then $\pi$ is an $\epsilon$-subgame perfect equilibrium of the Dec-POMDP, where $\epsilon=t_{max}T/e$.
\end{restatable}
Proof in the Appendix~\ref{app:proofs}.

\subsection{Optimal Grounded Policies}
\label{sec:optimal_grounded_policy}
In multi-agent cooperative settings with imperfect information, the optimal policy for an agent depends not only on assumptions about the future policy, but also on assumptions about other agents' policies -- i.e. what do those agents' prior actions say about their private information? Making wrong assumptions is a source of coordination failure and we thus may wish to consider grounded policies that avoid reasoning about partners' actions altogether. 

To do so, we define the \textit{grounded belief} $\mathcal{B}_G$ as a modified belief that conditions only on the observations but not the partner actions, i.e. $$\mathcal{B}_G(\tau|\tau^i) = \frac{P(\tau) \prod_t{  P( o_t^i | \tau ) } }{ \sum_{\tau'} P(\tau') \prod_t{  P( o_t^i | \tau' ) } }.$$ Then, we can define an \textit{optimal grounded policy} $\pi_G$ to be any policy that at each AOH  $\tau^i$ plays an action that maximizes expected reward assuming the state distribution at $\tau^i$ is drawn from the grounded belief, and assuming that $\pi_G$ is played thereafter.
\begin{restatable}{theorem}{thmgrounded}
\label{thm:grounded}
Application of OBL to any constant policy $\pi_0(a|\tau^i)=f(a)$ - or in fact any policy that only conditions on public state - yields an optimal grounded policy in the limit as $T \to 0$ .
\end{restatable}
Proof in the Appendix~\ref{app:proofs}.

\subsection{Iterated Application of OBL}
\label{sec:iteratedobl}
While an optimal grounded policy can be an effective baseline in some settings, clearly there are many settings where reasoning about the partner's behavior improves performance. In humans this ability is commonly described as \emph{theory of mind} and a variety of experimental studies \cite{camerer2003,agranov2012449,kleiman2016coordinate} have shown that humans typically carry out a limited number of such reasoning steps, depending on the exact situation, even in zero-shot settings.

Iterating OBL is simple. When the OBL operator is applied to a random policy $\pi_0$, it computes an optimal grounded policy $\pi_1$, which we also refer to as OBL ``level 1''. We may compute OBL level 2 similarly by using OBL level 1 as $\pi_0$ and compute OBL level 3 by replacing $\pi_0$ with OBL level 2, and so on.

Level 2 acts optimally while using beliefs induced by level 1 to interpret partners' actions. Different from OBL level 1, a level 2 agent will assume that other agents will take positive-utility actions instead of purely random actions and share useful factual information to others. For example, if it sees another agent avoid an action that would likely be beneficial, it will infer that the other agent probably has private information that the action is not so beneficial.

Since all independently trained agents use the same belief (modulo convergence errors), this training procedure is helpful in ZSC, unlike self-play where agents may learn arbitrary beliefs and signaling. Iterated OBL beliefs also appear to be a better match for human behavior than any prior work. We also show in Section~\ref{sec:result} that while OBL level 1 already cooperates better with a human-like policy than any prior method, level 2 and higher cooperate even better with it and other tested agents.

The iterated application of OBL allows us to directly control for the number of cognitive reasoning loops to be carried out by our agent. Different from the cognitive hierarchies framework, at each level the agent plays optimally assuming optimal game play in the future, rather than assuming the partners are fixed and play one level below. Therefore it is possible to optimize the overall \emph{quality} of the policy independently of the depth of cognitive reasoning in OBL but not in cognitive hierarchy because we cannot reach the optimal policy at each level.

\begin{figure*}[t]
\centering
\includegraphics[width=0.45\linewidth]{figs/fig-iql.pdf}
\includegraphics[width=0.45\linewidth]{figs/fig-obl.pdf}
\caption{\small Comparison of independent Q learning (left) with LB-OBL (right). Superscripts $i$ and $-i$ mean player $i$ and its partner $-i$. In LB-OBL, the target value for an action $a^i_t$ is computed by first sampling a state from a belief model $\hat{B}_{\pi_0}$ that assumes $\pi_0$ has been played, and then simulating play (blue) from that state to $i$'s next turn using the current policy $\pi_1$.}
\label{fig:iql-fig}
\end{figure*}

\subsection{Simultaneous-action vs turn-based environments}
\label{sec:turn_based}

So far we have restricted ourselves to \textit{turn-based} games in which a single agent acts at each state. If two players acted simultaneously, then OBL couldn't guarantee convergence to a unique policy (Thm \ref{thm:unique}); for example, consider a one-step coordination game with payoff matrix $[[1, 0], [0, 1]]$ (which is a single-state, single-step Dec-POMDP). OBL is identical to self-play in a one step game, and could thus converge to either equilibrium.

This may seem like a strong restriction of OBL to turn-based games. However, any simultaneous-action game can be converted to an equivalent stage game by ordering the simultaneous actions into separate timesteps and only letting players observe the actions of others once the timesteps corresponding to the simultaneous actions have completed. An OBL policy can be generated for this equivalent turn-based game. The ordering of players' actions solves the equilibrium selection problem, because the player who acts second assumes that the other player acted according to $\pi_0$. The player order we use to convert the simultaneous-action game to a turn-based one may change the OBL policy.

\subsection{Algorithms for Off-Belief Learning}
\label{sec:algo-obl}

Equation \ref{eq:pi1} immediately suggests a simple algorithm for computing an OBL policy in small tabular environments: compute $\mathcal{B}_{\pi_0}(\tau^i)$ for each AOH, and then compute $Q^{\pi_0\to\pi_1}(\tau^i)$ for each AOH in `backwards' topological order. However, such approaches for POMDPs are intractable for all but the smallest size problems.

In order to apply value iteration methods, we can write the Bellman equation for $Q^{\pi_0 \to \pi_1}$ for each agent $i$ as follows:
\begin{equation}
\begin{split}
    Q^{\pi_0\to \pi_1}(a_t|\tau_t^i) = \mathbb{E}_{\tau_t, \tau_{t+k}} \Big[ \sum_{t'=t}^{t+k-1} R(\tau_{t'}, a_{t'}) + \\
    \sum_{a_{t+k}} \pi_1(a_{t+k}|\tau_{t+k}^i) Q^{\pi_0\to \pi_1}(a_{t+k}|\tau_{t+k}^i)\Big]
\end{split}
\label{eq:ob_bellman}
\end{equation}
where $\tau_{t+k}$ is the next history at which player $i$ acts, $\tau_t \sim \mathcal{B}_{\pi_0}(\tau^i_t)$, and $\tau_{t+k} \sim (\mathcal{T},\pi_1)$ denotes that $\tau_{t+k}$ is sampled from the distribution of successor states where all (other) players play according to $\pi_1$. 

Now, how can we perform the Bellman iteration in Eq. \ref{eq:ob_bellman}? One approach, which we will denote \emph{Q-OBL}, is to perform independent $Q$-learning where each player uses $\pi_0$ as the exploration policy to play until time $t$ and use $\pi_1$ afterwards. We can then train $\pi_1$ at time $t$ but not other timesteps.
This would guarantee that $\tau_t \sim \mathcal{B}_{\pi_0}(\tau_t^i)$. However in order to compute the value for playing $\pi_1$ in the future, we must simulate the transition $\tau_{t+k}\sim(\mathcal{T},\pi_1)$ by having the other players playing according to their current $\pi_1$. This is guaranteed to converge to the OBL policy (in the tabular case) as long as $\forall \tau^i\ \forall a\ \pi_0(a|\tau^i) > 0$, by the same inductive argument used in Theorems \ref{thm:unique} and \ref{thm:policy_improvement}, since $Q$ and $\pi_1$ at each AOH only depend on $Q$ and $\pi_1$ at successor AOHs. 


\begin{figure*}[h]
\begin{center}
  \includegraphics[width=1.0\textwidth]{figs/results_toy_game.pdf}
      \vspace{-18pt}
  \caption{Cross-play matrix for each of the different algorithms in the toy communication game. 10 agents were independently trained with each method; each cell (i, j) of the matrix represents the average score of the $i$th and $j$th agent. OBL with a depth of 1 \textbf{(b)} achieves a higher score than CH \textbf{(e)} by learning optimal grounded communication, while avoiding arbitrary handshakes that lead to the poor XP performance of SP \textbf{(a)} or multi-level reasoning \textbf{(d, g)}.}
      \vspace{-8pt}
  \label{figs:simple_game_results}
\end{center}
\end{figure*}

In practice, this approach has a downside: the states reached by $\pi_1$ may be reached with very low probability under $\pi_0$,\footnote{The bound is $P(\tau|\pi_0)/P(\tau|\pi_1) \geq (\min_a P(a))^t$} so this procedure will have low sample efficiency. Instead, we first learn an \textit{approximate belief} $\hat{\mathcal{B}}_{\pi_0}$ that takes an AOH $\tau^i$ and can sample a trajectory from an approximation of $P(\tau|\tau^i,\pi_0)$. We compute this belief model following the procedure described in \citet{hu2021learned}. We then perform $Q$-learning with a modified rule for computing the target value for $Q(a|\tau_t^i)$: we re-sample a \textit{new} $\tau'$ from $\hat{B}_{\pi_0}(\tau_t^i)$, simulate a transition to $\tau_{t+1}^{i'}$ with other agents play their policy $\pi_1$, and use $\max_a Q(a|\tau_{t+1}^i)$ as the target value. We refer to this variant as \textit{Learned-belief OBL} (LB-OBL).

An illustration comparing the LB-OBL algorithm with independent Q-learning (IQL) is shown in Figure~\ref{fig:iql-fig}. In both cases, we assume a two player setting with shared Q-function and 2-step Q-learning, i.e. $k=2$ in Equation~\ref{eq:ob_bellman}. The active player at time $t$ is $i$ and the partner is $-i$. In IQL, each player simply observes the world at each time-step and takes turns to act. The target is computed on the actual trajectory with $r_t$, $r_{t+1}$ and $\max_{a}Q(a|\tau_{t+2})$ using target network $Q^{i}$. By contrast, LB-OBL involves only \emph{fictitious transitions}. For the active player $i$, we first decide the action $a_t^{i}$ given AOH $\tau_t^i$. However, in addition to applying $a_t^{i}$ to the actual environment state $S_t$, we also apply $a_t^{i}$ to a fictitious state $S^{\prime}_{t}$ sampled from the learned belief model $B_i$ and forward the fictitious environment to $S^{\prime}_{t+1}$. Note that the action $a_t^i$ can be directly applied to the fictitious state because the observation from the fictitious state is the same as that of the real state. Next, we need to evaluate the partner's policy to produce a fictitious action $a^{-i\prime}_{t+1}$ on $\tau^{-i\prime}_{t+1} = [\tau^{-i}; a^{i}_t, \Omega^{-i}(S^{\prime}_{t+1})]$. The learning target in LB-OBL is the sum of \emph{fictitious} rewards $r'_{t}$, $r'_{t+1}$ and the \emph{fictitious} bootstrapped value $\max_{a}Q(a|\tau'_{t+2})$. We also note that it is not possible to simply train a Q-learning agent that takes the \emph{grounded beliefs} as an input. The agents would still develop conventions and the \emph{grounded belief} would simply lose its semantic meaning. 
